\lesson{4}{Mar 18 2022 Fri (17:41:04)}{Simplify powers of $i$}{Unit 4}
\label{les_4_unit_4:simplify_powers_of_i_}

Let's say we're given the problem: $i^{40}$. We know that:

\begin{equation}
  \label{spt:powers_of_i}
  \begin{split}
    i^{1} &= \sqrt{-1} \\
    i^{2} &= -1 \\
    i^{3} &= -i \\
    i^{4} &= 1 \\
  \end{split}
.\end{equation}

To solve problems like this, you divide the power of $i$ by $4$, since that's
how many options you have, and the remainder, you raise $i$
to that power.

For this problem, you divide $40 \div 4 = 4 \qquad$. So, you get $i^{1} = i$.

\begin{example}
  \label{exm:simplfy_powers_of_i}

  Solve $i^{67}$. First, divide $67 \div 4 = 16 \qquad R = 3$, which gives you
  $i^{3} = -i$.
\end{example}

\newpage
